\chapter{Pneumococcal Disease Testing}

\section*{What Is This Test?}
Pneumococcal disease testing detects \textit{Streptococcus pneumoniae} (pneumococcus), a gram-positive, alpha-hemolytic, encapsulated diplococcus and the most common cause of community-acquired pneumonia, bacterial meningitis in adults, and occult bacteremia in children \cite{vanderpoll2009pathogenesis}. \textit{S. pneumoniae} has more than 100 serotypes defined by capsular polysaccharide; the 13-valent conjugate vaccine (PCV13) and 23-valent polysaccharide vaccine (PPSV23) cover the most clinically significant serotypes. Testing methods include: Gram stain and culture of clinical specimens (sputum, blood, CSF, pleural fluid); urinary antigen test (BinaxNOW \textit{S. pneumoniae}) detecting C-polysaccharide cell wall antigen common to all serotypes; and PCR detecting pneumococcal DNA targets (\textit{lytA}, \textit{ply}). The urinary antigen test has a sensitivity of 74--86\% for bacteremic pneumococcal pneumonia and 50--80\% for non-bacteremic pneumonia, with specificity >95\% \cite{sinclair2013systematic}. Blood cultures are positive in only 20--30\% of pneumococcal pneumonia cases. Antimicrobial susceptibility testing is essential due to increasing penicillin and ceftriaxone resistance.

\section*{Alternative Names}
Pneumococcus Testing; \textit{S. pneumoniae} Culture; Pneumococcal Urinary Antigen; BinaxNOW Pneumococcal Antigen; Strep Pneumoniae Susceptibility; Pneumococcal PCR

\section*{Principle}
Urinary antigen: Immunochromatographic membrane assay (BinaxNOW) detecting \textit{S. pneumoniae} C-polysaccharide (teichoic acid) cell wall antigen in urine. Antigen is detectable within 1--3 days of symptom onset. Sensitivity: 74--86\% (bacteremic pneumonia), 50--80\% (non-bacteremic pneumonia). Specificity: >95\%, though false positives occur in children with pneumococcal nasopharyngeal colonization (not performed in children due to low specificity). False positives after PCV13 vaccination are rare (antigen test is negative within 48 hours of vaccination in most individuals). Culture: Specimens inoculated on 5\% sheep blood agar; incubate at 35--37\textdegree{}C, 5\% CO$_2$. Colonies are alpha-hemolytic (green zone), mucoid (encapsulated), and characterized by a central depression (autolysis) in older colonies. Identification: Optochin susceptibility (zone $\geq$14 mm with 5 mcg disk) and bile solubility (lysis by sodium deoxycholate) \cite{werno2008medical}. Serotyping: Quellung reaction using type-specific antisera (reference laboratories). PCR: Real-time PCR targeting \textit{lytA} (autolysin) and \textit{ply} (pneumolysin) genes. Antimicrobial susceptibility: Penicillin MICs determined by broth microdilution or Etest; for meningitis, penicillin MIC $\geq$0.06 mcg/mL = resistant; for non-meningitis isolates, penicillin MIC $\leq$2 mcg/mL = susceptible (using new CLSI 2008 breakpoints) \cite{clsi2023performance}.

\section*{History}
\textit{S. pneumoniae} was first isolated independently by Louis Pasteur and George Sternberg in 1880--1881. Frederick Griffith's 1928 experiment demonstrating transformation of nonvirulent to virulent pneumococci (by DNA) was foundational to molecular biology. Oswald Avery, Colin MacLeod, and Maclyn McCarty in 1944 identified DNA as the transforming substance. The first pneumococcal polysaccharide vaccine was licensed in 1977 (14-valent), followed by PPSV23 in 1983. The 7-valent conjugate vaccine (PCV7) was introduced in 2000, dramatically reducing invasive pneumococcal disease (IPD). PCV13 replaced PCV7 in 2010, and PCV15 and PCV20 were licensed in 2021. The urinary antigen test (BinaxNOW) was FDA-cleared in 2003 and has become a standard component of the diagnostic evaluation for hospitalized CAP. Penicillin-nonsusceptible \textit{S. pneumoniae} emerged in the 1970s and became widespread by the 1990s; however, rates of penicillin resistance have declined since the introduction of conjugate vaccines.

\section*{Why This Test Is Ordered}
\begin{itemize}
  \item Diagnosis of the causative pathogen in patients hospitalized with community-acquired pneumonia, to allow targeted (narrowest-spectrum) antibiotic therapy and de-escalation from broad-spectrum empiric regimens
  \item Rapid diagnosis of invasive pneumococcal disease (bacteremia, meningitis) in critically ill patients where time to appropriate therapy impacts mortality
  \item Evaluation of patients with pyogenic meningitis (CSF pleocytosis with gram-positive diplococci on Gram stain) for definitive diagnosis and antimicrobial susceptibility
  \item Investigation of occult bacteremia in febrile children 3--36 months without a focus of infection (though incidence has declined >90\% since PCV introduction)
  \item Diagnosis of pneumococcal peritonitis (spontaneous bacterial peritonitis in patients with cirrhosis or nephrotic syndrome, or primary peritonitis in adolescent females)
  \item Public health surveillance for invasive pneumococcal disease, including serotype distribution for vaccine impact monitoring
  \item Antibiotic stewardship: identification of \textit{S. pneumoniae} with penicillin susceptibility allows safe de-escalation from vancomycin plus broad-spectrum beta-lactam to penicillin or amoxicillin for susceptible isolates
\end{itemize}

\section*{Primary vs Secondary Test}
The pneumococcal urinary antigen test is a primary diagnostic test recommended for hospitalized patients with CAP (IDSA/ATS 2019 guidelines) \cite{metlay2019diagnosis}. Blood cultures should be obtained concurrently. Sputum Gram stain and culture, when an adequate specimen can be obtained, are complementary primary tests. CSF culture and Gram stain are primary in suspected meningitis. PCR may be used as a secondary, confirmatory test when culture is negative.

\section*{How This Test Is Performed}
Urinary antigen: Random clean-catch urine specimen (3--5 mL) in a sterile cup; no preservatives; refrigerate if not tested immediately. Blood cultures: Two sets (aerobic and anaerobic bottles) from separate venipuncture sites before antibiotic administration; 20 mL of blood per set in adults. Sputum culture: Expectorated deep cough specimen after mouth rinse; evaluate Gram stain for specimen adequacy (<10 SECs/LPF, >25 WBCs/LPF). CSF: Collected by lumbar puncture in sterile tubes; Tube 1 for chemistry/glucose/protein, Tube 2 for microbiologic studies, Tube 3 for cell count/differential, Tube 4 for additional testing. Transport to laboratory within 15 minutes (CSF STAT) or within 2 hours (other specimens). Do not refrigerate CSF for culture as this reduces recovery of fastidious organisms.

\section*{What Preparation Is Needed}
No special preparation for urine, sputum, or blood collection. Antibiotics should be withheld until after cultures are obtained if the clinical situation permits, as even a single dose of an effective antibiotic can reduce blood culture sensitivity by up to 50\%.

\section*{Contraindications and Precautions}
The pneumococcal urinary antigen test is NOT recommended in children, as asymptomatic nasopharyngeal pneumococcal colonization (present in 20--60\% of healthy children) may cause false-positive results. It should be used with caution in patients with recent (within 5 days) PCV13 vaccination, though false positives are uncommon. A negative urinary antigen does NOT exclude pneumococcal pneumonia, especially non-bacteremic disease. In patients with COPD or chronic bronchial colonization, a positive sputum culture may represent colonization rather than acute infection, and correlation with Gram stain, clinical findings, and blood culture results is essential.

\section*{Risks}
Risk from urinary antigen testing is negligible. Risks from blood culture, lumbar puncture, or sputum induction are related to the collection procedure, not the laboratory analysis. The major clinical risk is failure to identify invasive pneumococcal disease rapidly, as mortality increases significantly with delays in appropriate antibiotic therapy (mortality 20--30\% for pneumococcal bacteremia in adults).

\section*{What Happens After the Test}
Urinary antigen results: 15--20 minutes (BinaxNOW). Sputum Gram stain: 1--2 hours. Blood culture: preliminary positive at 12--24 hours (Gram stain); final identification and sensitivity at 48--72 hours. CSF culture and Gram stain: Gram stain within 1 hour STAT; culture preliminary at 24 hours. Positive blood or CSF cultures with gram-positive cocci suggesting pneumococcus are immediately communicated to the provider as a critical result. Empiric therapy for suspected pneumococcal meningitis should include vancomycin plus ceftriaxone until penicillin/ceftriaxone susceptibility is confirmed, due to the risk of resistant isolates.

\section*{Understanding Results}

\subsection*{Below Normal}
\begin{itemize}
  \item \textbf{Urinary antigen negative:} Does not exclude pneumococcal infection. 14--26\% of bacteremic pneumococcal pneumonia and up to 50\% of non-bacteremic pneumococcal pneumonia may have negative urinary antigen. Causes: early infection (antigen not yet detectable); non-pneumococcal etiology; low bacterial burden; intermittent antigen excretion. The test's sensitivity is higher in bacteremic than non-bacteremic disease
  \item \textbf{Blood cultures negative:} Only 20--30\% of pneumococcal pneumonia is bacteremic. A negative blood culture has a negative predictive value of 70--80\% for excluding pneumococcal etiology
  \item \textbf{Sputum culture negative or normal flora:} May represent an inadequate specimen, prior antibiotic therapy, or non-pneumococcal etiology
\end{itemize}

\subsection*{Above Normal}
\begin{itemize}
  \item \textbf{Urinary antigen positive:} Indicates \textit{S. pneumoniae} C-polysaccharide antigen in urine. In an adult patient with pneumonia, this is highly suggestive of pneumococcal pneumonia (PPV >95\%). False positives may occur in patients receiving PCV13 vaccination within the preceding 48 hours and in patients with pneumococcal colonization and concurrent non-pneumococcal pneumonia. The test does not distinguish among serotypes
  \item \textbf{Blood culture positive for \textit{S. pneumoniae}:} Definitive diagnosis of invasive pneumococcal disease. All isolates should undergo antimicrobial susceptibility testing for penicillin (meningitis and non-meningitis breakpoints), ceftriaxone (meningitis breakpoint: MIC $\leq$0.5 mcg/mL), erythromycin, clindamycin, trimethoprim-sulfamethoxazole, vancomycin, and levofloxacin. Penicillin-nonsusceptible rates are higher for isolates associated with meningitis, bacteremia, and otitis media
  \item \textbf{Sputum culture predominant \textit{S. pneumoniae} + compatible Gram stain:} When correlated with a compatible Gram stain (many WBC, >25/LPF; gram-positive lancet-shaped diplococci; few SEC, <10/LPF), indicative of pneumococcal pneumonia
  \item \textbf{CSF positive (Gram stain + culture + antigen):} Diagnostic of pneumococcal meningitis. A negative CSF Gram stain does not exclude the diagnosis; empiric therapy should not be delayed
\end{itemize}

\section*{Normal Results}
\textbf{Negative} urinary antigen for \textit{S. pneumoniae}. Blood culture: \textbf{No growth at 5 days}. Sputum culture: \textbf{Normal respiratory flora} or \textbf{No \textit{S. pneumoniae} isolated}. CSF: \textbf{No organisms on Gram stain, no growth at 5 days}.

\section*{Follow-up Tests}
\begin{itemize}
  \item \textbf{Antimicrobial susceptibility testing (AST):} For all \textit{S. pneumoniae} isolates from normally sterile sites (blood, CSF, pleural fluid) and from sputum in cases of clinical pneumonia; includes penicillin/oxacillin disk diffusion screen, confirmation by broth microdilution MIC, ceftriaxone/cefotaxime MIC (if meningitis), and susceptibility to macrolides, fluoroquinolones, clindamycin, trimethoprim-sulfamethoxazole, vancomycin
  \item \textbf{Pneumococcal serotyping:} Quellung reaction or PCR-based serotyping of isolates from normally sterile sites for public health surveillance and vaccine impact monitoring; performed at reference/public health laboratories
  \item \textit{\textbf{S. pneumoniae}} \textbf{PCR on CSF or pleural fluid:} When cultures are negative but clinical suspicion remains high (administered antibiotics prior to CSF collection, partially treated meningitis). PCR targeting \textit{lytA} provides high sensitivity and specificity
  \item \textbf{Chest radiograph (PA and lateral, or CT chest):} To evaluate the extent of pneumonia (lobar consolidation, multilobar disease, parapneumonic effusion, empyema, cavitation)
\end{itemize}

\section*{References}
\bibliographystyle{unsrtnat}
\bibliography{references}

\clearpage
\section*{Regulatory Status and Authorization Date}
Regulatory status applies to the specific assay, instrument, manufacturer, and intended use rather than to the test name alone. No single approval date can be assigned to this test category. For a commercial assay, verify the current product in the relevant regulator's database: the U.S. Food and Drug Administration (FDA) clearance, approval, or authorization; India's Central Drugs Standard Control Organisation (CDSCO) licence or permission; the European Union In Vitro Diagnostic Regulation (EU IVDR) CE certificate issued through the applicable conformity-assessment route; or the United Kingdom Medicines and Healthcare products Regulatory Agency (MHRA) registration and UKCA/CE status. Laboratory-developed tests may not have an individual FDA, CDSCO, EU IVDR, or MHRA approval date.
\clearpage
